\maketitle
\thispagestyle{plain}
\begin{abstract}
\begingroup
\linespread{1.14}\selectfont

光伏出力、居民负荷和外网电价都具有明显的时变性，储能又把相邻时段联结起来，微网购电由此形成带状态约束、预测误差和计划偏差结算的连续决策问题。围绕题目给出的四种信息条件，本文始终保持微网能量平衡与储能状态递推不变，仅随问题推进改变计划时可用的信息、允许的调整方式及相应费用，从而形成统一的日前—执行—滚动调度框架。

\textbf{对于问题一}，在典型日电价、负载和光伏预测均已知时建立确定性线性规划，以全天购电费最小为目标，联合决定外网购电与储能充放电，并固定日初、日末储电量为 $6000$ kWh。求得全天计划购电量为 \QoneEnergy\ kWh，购电费为 \QoneCost\ 元；与禁用储能时的 \QoneNoStorage\ 元相比，费用降低 \QoneSave\ 元，降幅为 \QoneSavePct\%。结果表明，储能主要通过低价时段充电、较高电价时段放电以及吸收中午光伏富余实现跨时段能量转移（详见第~\pageref{tab:q1-result} 页表~\ref{tab:q1-result}）。

\textbf{对于问题二}，每天 $0{:}00$ 制定的购电计划必须先于当天实际源荷实现。为避免把未来真实数据提前用于计划，主模型只利用此前历史数据预测当天负载和光伏，并将计划购电量在日内锁定。针对紧急购电价格为正常电价 $5$ 倍这一非对称代价，利用单时段边际损失得到光伏低分位风险修正的解析依据，再在历史窗口中滚动选择预测器与安全边际；执行阶段允许储能根据已实现偏差实时缓冲，仍不能覆盖的部分才紧急购电。2月至12月总费用为 \QtwoTotal\ 元，紧急购电量为 \QtwoEmergency\ kWh。对照实验表明，负荷信息边界、光伏风险修正和储能实时反馈都会显著影响全年成本（详见第~\pageref{tab:q2-ablation} 页表~\ref{tab:q2-ablation}）。

\textbf{对于问题三}，在 $0{:}00$ 先形成全天基准计划，在 $6{:}00$、$12{:}00$、$18{:}00$ 获得新预报后，只对尚未执行的时段重新优化，并以当时实测储电量衔接前后两个滚动阶段。结算时，未被调整的部分按正常电价计费，少购的部分按电价的 $50\%$ 支付违约费用，超出计划的部分按电价的 $150\%$ 计费；本文据此构造主模型，并将“计划购电费与调整费用并列相加”的另一种理解作为对照。仅在 $0{:}00$ 决策时，全年费用为 \QthreeSzero\ 元；依次加入三个日内更新节点后，费用降至 \QthreeSone、\QthreeStwo、\QthreeSthree\ 元，累计节省 \QthreeVtotal\ 元。配对对照表明，节点收益由已实现负荷信息与新光伏预报共同形成，本文将其解释为决策更新的综合价值，其中光伏刷新的增量单独报告（详见第~\pageref{tab:q3-ablation} 页表~\ref{tab:q3-ablation}）。

\textbf{对于问题四}，保持问题二、三的预测与滚动结构，将固定日电价替换为附件4给出的逐日电价重新优化。按题面直接使用附件4价格时，Q4-2 与 Q4-3 的全年费用分别为 \QfourTwoMain\ 元和 \QfourThreeMain\ 元，滚动更新节省 \QfourRollPerfect\ 元。以历史价格预测替代未来完整价格作为现实性扩展后，价格信息受限会增加费用，但其影响小于源荷信息误差，且不改变日内滚动更新能够降低总成本的结论（详见第~\pageref{tab:q4-result} 页表~\ref{tab:q4-result}）。

为保证结果满足物理可行性与时间可执行性，全文分别检验能量平衡、储能状态递推、容量和功率边界、计划持续性、费用回代以及信息因果性，并对终端储电量和效率取值进行敏感性分析。结果说明：储能提供了跨时段调节能力，对预测信息、实际观测和计划结算边界的严格区分，是避免微网调度结果过度乐观的关键。

\keywords{微网经济调度；储能优化；因果预测；风险购电；滚动优化}
\endgroup
\end{abstract}
